\unitopen{48}
\begin{candos}
\item przejść test końcowy B2 z minimum 80\% poprawnych odpowiedzi;
\item przeprowadzić trzydziestominutową rozmowę zawodową po niemiecku z AI lub korepetytorem i ocenić ją według kryteriów egzaminu;
\item zdecydować, czy zdawać egzamin certyfikujący i zaplanować go;
\item ułożyć plan na kolejny rok (utrzymanie, C1, dialekt, pierwsze miesiące w pracy).
\end{candos}

\section{Słownictwo}
\vocabtable{48}

\section{Format egzaminu B2}
Egzaminy takie jak Goethe-Zertifikat B2, telc Deutsch B2 (także wariant zawodowy) i ÖSD ZB2 mają moduły czytania, słuchania, pisania i mówienia.
Mówienie sprawdza prezentację tematu, rozmowę i reakcję na pytania; pisanie sprawdza np.\ list formalny i tekst argumentacyjny
\verify{(sprawdź aktualny format, czas i punktację w arkuszu modelowym wystawcy egzaminu)}. Do tej pory przećwiczyłeś wszystkie sprawności: teraz weryfikujesz je jednym
pełnym arkuszem w warunkach egzaminu.

\begin{center}\small
\begin{tabular}{@{}p{0.2\linewidth}p{0.72\linewidth}@{}}
\toprule
\textsf{\textbf{Moduł}} & \textsf{\textbf{Co oceniają egzaminatorzy}}\\
\midrule
\de{Lesen} & rozumienie sedna i szczegółów, wnioskowanie, funkcje tekstu\\
\de{Hören} & rozumienie wiadomości, rozmów, wykładów w normalnym tempie\\
\de{Schreiben} & spójność, struktura argumentu, rejestr, poprawność, zakres słownictwa\\
\de{Sprechen} & płynność, treść, gramatyka, słownictwo, wymowa, interakcja\\
\bottomrule
\end{tabular}
\end{center}

\begin{zasada}[Próg 80\% i uczciwa ocena]
Jeśli test końcowy poniżej 80\% (8 z 10), wróć do jednostek 37, 39, 41. Jeśli zdajesz egzamin próbny w czterech modułach, każdy powinien dać co najmniej
60--70\% punktów \verify{(próg zaliczenia zależy od egzaminu)}. Szukaj słabego modułu, nie średniej: jeden moduł pod progiem oznacza dodatkowy tydzień.
\end{zasada}

\begin{falle}
Po roku intensywnej nauki łatwo pomylić zrozumiałość z poziomem. Błąd: uznanie B2 po jednym dobrym nagraniu. Drugi: pomijanie pisania, bo „i tak
mówię”. Trzeci: zaniedbanie słownictwa po zakończeniu kursu: bez powtórek Anki zasób maleje szybciej niż rośnie.
\end{falle}

\begin{szwajcaria}
Praktyka szwajcarska: certyfikat B2 nie jest zwykle wymagany od pracownika IT z UE, ale dobra niemiecka komunikacja zwiększa Twoje szanse w
firmach z klientem lokalnym, w administracji i w pracy z zespołami niemieckojęzycznymi. Najlepszy dowód to rozmowa kwalifikacyjna po niemiecku, a drugi to
referencje. Jeśli planujesz pobyt długoterminowy, sprawdź w SEM i w kantonie, czy poziom językowy jest wymagany dla Twojej sytuacji
\verify{(warunki zależą od statusu i kantonu)}.
\end{szwajcaria}

\begin{praca}
Rozmowa 30 minut: 5 minut small talk, 10 minut projekt (prezentacja), 10 minut pytania techniczne i zespołowe, 5 minut negocjacje. Nagraj całość,
oceń według czterech kryteriów i wyślij do drugiej osoby. Dopiero ocena z zewnątrz mówi coś o poziomie.
\end{praca}

\section{Dialogi}
\textbf{Dialog 1 --- rozmowa o postępie}
\begin{dialog}
\turn{A}{Vor einem Jahr konnten Sie kaum etwas sagen. Wie fühlen Sie sich heute?}{Rok temu prawie nic Pan nie mówił. Jak się Pan czuje dziś?}
\turn{B}{Ich verstehe viel mehr, und ich kann in Meetings meine Meinung äußern. Schriftlich möchte ich aber noch sicherer werden.}{Rozumiem dużo więcej i potrafię wyrazić opinię na spotkaniach. Na piśmie chcę jednak być jeszcze pewniejszy.}
\turn{A}{Was ist Ihr Plan für das nächste Jahr?}{Jaki jest Pana plan na następny rok?}
\turn{B}{Ich werde regelmäßig lesen und schreiben, die Prüfung ablegen und mich später auf C1 vorbereiten.}{Będę regularnie czytać i pisać, zdam egzamin, a później przygotuję się do C1.}
\end{dialog}

\section{Test końcowy B2}
Próg: 80\% (8 z 10).

\begin{cwiczenia}
\exitem{Nominalizacja: \de{Wir prüfen den Antrag.} (w stylu formalnym)}{\de{Die Prüfung des Antrags erfolgt (durch uns).}}
\exitem{Określenie imiesłowowe: \de{die Daten, die im Test verwendet wurden}.}{\de{die im Test verwendeten Daten}}
\exitem{Mowa zależna: \de{Er sagte: „Ich bin krank.”}}{\de{Er sagte, er sei krank.}}
\exitem{Passiv stanu: \de{Der Test ist \dots{} (abschließen)}.}{\de{Der Test ist abgeschlossen.}}
\exitem{Połącz (\de{indem}): \de{Wir senken die Latenz. Wir führen einen Cache ein.}}{\de{Wir senken die Latenz, indem wir einen Cache einführen.}}
\exitem{\de{Je \dots{} desto}: \de{viele Tests / wenige Fehler}.}{\de{Je mehr Tests wir haben, desto weniger Fehler gibt es.}}
\exitem{Opisz wykres: spadek z 80 na 50 (\de{von \dots{} auf}).}{\de{Der Wert ist von 80 auf 50 gesunken.}}
\exitem{Uprzejma prośba o przesunięcie terminu w mailu.}{\de{Wäre es möglich, den Termin auf Donnerstag zu verschieben?}}
\exitem{Łagodnie wyraź sprzeciw na spotkaniu.}{\de{Ich sehe das etwas anders, weil \dots}}
\exitem{Przetłumacz: „Decyzję podjęliśmy, żeby zmniejszyć koszty, mimo że wymaga ona większego nakładu pracy.”}{\de{Wir haben die Entscheidung getroffen, um die Kosten zu senken, obwohl sie mehr Aufwand erfordert.}}
\end{cwiczenia}

\section{Trening słownictwa (generowany)}
\drill{48}

\weekend{\textbf{Zadanie końcowe B2.} Przeprowadź i nagraj trzydziestominutową symulację rozmowy zawodowej, napisz formalny
e-mail, ADR i tekst argumentacyjny (po 250 słów). Porównaj z nagraniami z A1, A2 i B1. Zapisz, co się zmieniło, i zaplanuj następny rok.}

\begin{uwaga}[Tydzień buforowy B2 i dalej]
Ostatni tydzień buforowy to egzamin próbny, powtórka Anki i plan: utrzymanie (codzienne czytanie i słuchanie), egzamin certyfikujący,
dialekt (rozumienie) i język w pracy w pierwszych miesiącach. Nauka nie kończy się na B2: do swobody zawodowej (C1) potrzeba kolejnych
setek godzin, ale masz już fundament, żeby się uczyć \emph{w pracy}.
\end{uwaga}
